\section{Discussion}
\label{sec:discussion}

\paragraph{What the theory buys a practitioner.}
The main output is not a fixed source fraction or a zero-mixing derivative test. It is a way to distinguish single-valley, multivalley, boundary, multiple-optimum, and flat regimes before choosing a mixing weight. Exact root isolation can certify global candidates for fixed parameters; finite-sample corrections then determine local shifts and select among population-optimal branches. Under matched spectra and fixed regularization, the gain expansion and first-order exchange approximation quantify data value, while the quadratic surrogate provides an interpretable dosage rule whose accuracy requires separate evaluation.

\paragraph{Limitations.}
The analysis uses fixed dimension, Gaussian covariates, a shared spectral basis, and fixed regularization. The log-convex classification is restricted to zero mismatch; the closed-form branch diagram currently covers a specific two-coordinate family, not the entire parameter space. General $R_1/R_n$ shape theory, higher-order rates at degenerate or multiple optima, tuned regularization, the proportional regime, and robot-policy validation remain open. The surrogate dosage rule omits Gram-fluctuation terms of the same asymptotic order as its retained noise terms and is not a complete full-ridge dosage guarantee.

\paragraph{Outlook.}
The next priorities are a broader classification of $R_1$ and $R_n$, higher-order degenerate perturbation theory, tuned-regularization baselines, and evaluation with independently estimated parameters on native cross-embodiment tasks. Such experiments should report both open-loop risk and closed-loop performance.

\paragraph{Conclusion.}
We presented a global analysis of full-ridge auxiliary-data mixing. Actual-risk expansions connect population geometry to finite-sample optima. Zero-mismatch log-convexity, minimum-dimension multimodality, exact algebraic certificates, and first-order branch selection replace a single assumed risk shape with a hierarchy of global structures and their consequences. The claims remain limited to the stated model and evidence hierarchy.
